\documentclass[12pt]{article}
\usepackage[utf8]{inputenc}
\usepackage{chemfig}
\usepackage[T1]{fontenc} % Ensures proper font encoding for copy-paste/accents
\usepackage{microtype}    % Improves justification and reduces hyphenation
\usepackage{amsmath, amssymb, amsthm}
\usepackage[letterpaper, margin=0.8in, heightrounded]{geometry}
\usepackage[skip=0.8em, indent=0pt]{parskip} % Modern way to handle paragraph spacing
\usepackage{setspace}
\setstretch{0.8} % Safer than \renewcommand{\baselinestretch}
\usepackage{graphicx}
\usepackage{float}
\usepackage{booktabs} % Essential for clean, professional-looking tables
\graphicspath{{images/}}
\usepackage{xcolor}
\usetikzlibrary{arrows.meta, matrix} % <-- Add 'matrix' here!
\usepackage{hyperref}
\hypersetup{
    colorlinks=true,
    linkcolor=blue,
    filecolor=magenta,      
    urlcolor=cyan,
    citecolor=green,
    pdftitle={My LaTeX Document},
    pdfpagemode=FullScreen,
}


\title{VESPR SHAPES ASSIGNMENT}
\author{Sarv Patel}
\date{\today}

\begin{document}
\maketitle

\section*{\normalsize Lewis \& Line Diagram}

\begin{figure}[htbp]
  \centering
  \chemfig[bond style={line width=0.8pt}, atom sep=2.3em]{**6(-(-N^{\oplus}(-[-3]\charge{180=:,270=..,0=:}{O}^{\text{ }\ominus})(=[-1]\charge{180=:,0=:}{O}))
  -
  -(-N^{\oplus}(-[-1]\charge{180=:,270=..,0=:}{O}^{\text{ }\ominus})(=[2]\charge{180=:,0=:}{O}))
  -(-CH_3)
  -(-N^{\oplus}(-[-3]\charge{180=:,270=..,0=:}{O}^{\text{ }\ominus})(=[2]\charge{180=:,0=:}{O}))
  -)}
  \quad
  \chemfig[bond style={line width=0.8pt}, atom sep=2.3em]{**6(-(-N^{\oplus}(-[-3]\charge{180=:,270=..,0=:}{O}^{\text{ }\ominus})(=[-1]\charge{180=:,0=:}{O}))
  -
  -(-N^{\oplus}(-[2]\charge{180=:,270=..,0=:}{O}^{\text{ }\ominus})(=[-1]\charge{180=:,0=:}{O}))
  -(-CH_3)
  -(-N^{\oplus}(-[-3]\charge{180=:,270=..,0=:}{O}^{\text{ }\ominus})(=[2]\charge{180=:,0=:}{O}))
  -)}
  \quad
  \chemfig[bond style={line width=0.8pt}, atom sep=2.3em]{**6(-(-N^{\oplus}(-[-1]\charge{180=:,270=..,0=:}{O}^{\text{ }\ominus})(=[-3]\charge{180=:,0=:}{O}))
  -
  -(-N^{\oplus}(-[2]\charge{180=:,270=..,0=:}{O}^{\text{ }\ominus})(=[-1]\charge{180=:,0=:}{O}))
  -(-CH_3)
  -(-N^{\oplus}(-[-3]\charge{180=:,270=..,0=:}{O}^{\text{ }\ominus})(=[2]\charge{180=:,0=:}{O}))
  -)}
  \caption{Trinitrotoluene (TNT) Structures}
  \label{fig:TNT}
\end{figure}


\section*{\small Physical Properties of TNT}


\section*{\small Structure \& Intermolecular Forces}

\section*{\small Comparision to a Similar Compound}

\begin{figure}[htbp]
  \centering
  \setchemfig{atom sep=2.5em}
  
    \label{fig:RDX_lewis}
    \chemfig[bond style={line width=0.8pt}]{*6(-N(-[-2]N(=[-1]O)(=[-3]O))--N(-N(=[-1]O)(=[2]O))--N(-N(=[2]O)(=[-3]O))-)}
    \quad
    \chemfig[bond style={line width=1pt}]{*6(C?(-H)=C(-[:-90]N(=[:210]O)(-[:-30]O))-C(-[:-30]H)=C(-[:30]N(=[:90]O)(-[:-30]O))-C(-[:90]C(<:H)(<[:30]H)(-[:150]H))=C?(-[:150]N(=[:90]O)(-[:210]O)))}

  \caption{RDX Vs. TNT}
  \label{fig:RDX}
\end{figure}

\begin{figure}[htbp]
  \centering
\end{figure}

\subsection*{\small Explanation of Difference: }

\end{document}
